\documentclass[11pt,twocolumn]{article}

% ─── Packages ─────────────────────────────────────────────────────────────────
\usepackage[utf8]{inputenc}
\usepackage[T1]{fontenc}
\usepackage{times}
\usepackage[margin=1in]{geometry}
\usepackage{graphicx}
\usepackage{amsmath,amssymb}
\usepackage{booktabs}
\usepackage{multirow}
\usepackage{hyperref}
\usepackage{xcolor}
\usepackage{algorithm}
\usepackage{algpseudocode}
\usepackage{caption}
\usepackage{subcaption}
\usepackage{natbib}
\usepackage{float}
\usepackage{enumitem}
\usepackage{url}

\hypersetup{
    colorlinks=true,
    linkcolor=blue!60!black,
    citecolor=blue!60!black,
    urlcolor=blue!60!black,
}

% ─── Title ────────────────────────────────────────────────────────────────────
\title{%
    \textbf{Machine Learning Detection of Pathogenic Short Tandem Repeat
    Expansions in Long-Read Sequencing Data}\\[6pt]
    \large{A Multi-Scale Windowed Classification Approach with\\
    Motif Decomposition and Clinical Annotation}
}

\author{
    Brandon Tang
    \\[4pt]
    Department of Computer Science, Simon Fraser University, Vancouver, Canada
}

\date{March 2026}

\begin{document}
\maketitle

% ══════════════════════════════════════════════════════════════════════════════
\begin{abstract}
Short tandem repeats (STRs) are genomic regions where a 1--6 base pair
motif is repeated in tandem.  Pathogenic expansions of specific STR loci
cause over 50 neurological and neuromuscular diseases, including
Huntington disease (HTT), Fragile X syndrome (FMR1), Friedreich ataxia
(FXN), and spinocerebellar ataxias (ATXN10, FGF14).  Long-read Oxford
Nanopore Technology (ONT) sequencing can span full repeat expansions
but introduces the challenge of detecting STR regions embedded within
reads exceeding 40\,kbp.  We present an open-source pipeline that
combines (i)~gradient-boosted classification with 47~engineered
features inspired by TRMotifAnnotator motif decomposition and STRling
$k$-mer counting, (ii)~a multi-scale sliding window strategy for
long-read STR localisation, and (iii)~REDatlas-derived clinical
annotation of detected loci.  On a balanced dataset of 5{,}986 sequences
from NCBI dbVar, our XGBoost classifier achieves 99.4\% F1 with
0.999~ROC-AUC under 5-fold stratified cross-validation.  We validate the
full pipeline on ONT haplotagged BAM files from the 1000 Genomes Project,
successfully detecting a pathogenic ATTCT$\times$983 expansion in the
ATXN10 locus of sample HG01122 and correctly classifying it as
spinocerebellar ataxia type~10 (SCA10) with ``full mutation'' status.
All code, models, and test data are publicly available.
\end{abstract}

\noindent\textbf{Keywords:} short tandem repeats, repeat expansion
disorders, long-read sequencing, machine learning, XGBoost,
motif decomposition, clinical genomics

% ══════════════════════════════════════════════════════════════════════════════
\section{Introduction}\label{sec:introduction}

Short tandem repeats (STRs), also called microsatellites, are genomic
loci where a short motif of 1--6 base pairs is repeated consecutively.
STRs constitute approximately 3\% of the human genome and are among
the most mutable loci, with mutation rates $10^3$--$10^5$ times higher
than point mutations~\citep{gymrek2017}.  While most STR variation is
benign, pathogenic expansions beyond locus-specific thresholds cause
a growing list of over 50 neurological and neuromuscular
diseases~\citep{depienne2021}.

The clinical significance of STR expansions was first established with
Fragile X syndrome (FMR1; CGG$\times$200+) and Huntington disease
(HTT; CAG$\times$40+)~\citep{fu1991,macdonald1993}.  More recently,
biallelic AAGGG expansions in RFC1 (CANVAS), pentanucleotide ATTCT
expansions in ATXN10 (SCA10), and intronic GAA expansions in FGF14
(SCA27B) have been characterised~\citep{cortese2019,matsuura2000,pellerin2023}.
Accurate genotyping of these loci is essential for clinical diagnosis
and genetic counselling.

\subsection{Challenges in STR Detection}

Short-read whole-genome sequencing (Illumina, $\sim$150\,bp reads) cannot
span expansions exceeding a few hundred base pairs, leading to systematic
underestimation of long alleles~\citep{tankard2018}.  Long-read
technologies---Oxford Nanopore Technology (ONT) and PacBio HiFi---can
produce reads exceeding 100\,kbp, fully spanning even massive expansions
such as DMPK CTG$\times$1000+ (myotonic dystrophy type 1).  However,
long reads introduce new computational challenges:

\begin{enumerate}[nosep]
    \item \textbf{Repeat dilution}: An STR expansion of $\sim$5\,kbp
          embedded in a 40\,kbp read constitutes only $\sim$12.5\% of the
          total sequence.  Whole-read features (entropy, $k$-mer diversity)
          are dominated by flanking non-repeat sequence, masking the
          repeat signal.
    \item \textbf{Error rates}: ONT reads have $\sim$5--10\% per-base
          error rates (R9.4.1 chemistry), introducing substitutions and
          indels within repeat tracts that complicate motif identification.
    \item \textbf{Motif ambiguity}: At high error rates, short motifs
          (e.g., mononucleotide runs) can spuriously match portions
          of longer true motifs (e.g., ``T'' matching 40\% of ``ATTCT'').
\end{enumerate}

\subsection{Related Work}

Several tools address STR genotyping from long reads.
\textbf{STRling}~\citep{dashnow2022} uses non-overlapping $k$-mer
counting with canonical rotation to detect STR-containing reads,
followed by locus-level clustering and genotyping.
\textbf{RExPRT}~\citep{dolzhenko2024} employs a random forest
classifier with 56~features encompassing genomic architecture,
repeat characteristics, and population-level information to predict
pathogenicity of tandem repeat expansions.
\textbf{HMMSTR}~\citep{steely2025} applies hidden Markov models to
long-read sequences, using Viterbi decoding for single-base-resolution
motif copy counting even in the presence of sequencing errors.
\textbf{TRMotifAnnotator}~\citep{wf-trs2025} provides motif
decomposition that separates canonical and variant (interrupted) repeat
copies, enabling purity assessment.
\textbf{REDatlas}~\citep{rajanbabu2025} catalogues 66
disease-associated TR loci across 2{,}530 haplotypes from the 1000
Genomes ONT dataset, providing population allele frequencies and
clinical thresholds.

\subsection{Contributions}

We present a complete pipeline that integrates ideas from the above
tools into a single system:

\begin{enumerate}[nosep]
    \item A 47-feature XGBoost classifier combining TRMotifAnnotator-inspired
          motif decomposition, STRling-inspired $k$-mer counting, and
          traditional sequence composition features.
    \item A multi-scale sliding window strategy that overcomes repeat
          dilution in long reads via coarse (5\,kbp) then refined (1\,kbp)
          window passes.
    \item Automated clinical annotation via an integrated REDatlas
          module mapping detected loci to 21~disease-associated TR loci
          with pathogenicity classification.
    \item Validation on real ONT data from the 1000 Genomes Project,
          demonstrating correct detection of a pathogenic ATXN10
          expansion.
\end{enumerate}

% ══════════════════════════════════════════════════════════════════════════════
\section{Methods}\label{sec:methods}

\subsection{Training Data Acquisition}\label{sec:data}

Training data was sourced from NCBI dbVar~\citep{lappalainen2013} via
the NCBI E-utilities API.  We queried for structural variants of type
``short tandem repeat'' or ``microsatellite'' with a minimum length of
6\,bp.  Reference genome sequences (GRCh37) were retrieved for each
variant locus using the NCBI efetch interface.

\textbf{Positive class} ($n = 2{,}993$): STR sequences from dbVar
with clinical significance annotations including pathogenic, likely
pathogenic, and benign STR variants spanning unit lengths 1--6\,bp.

\textbf{Negative class} ($n = 2{,}993$): Non-STR genomic sequences
sampled from random genomic regions with matched length distributions,
filtered to exclude any overlap with known STR loci.

The balanced dataset ($N = 5{,}986$) was stratified by chromosome and
repeat unit length to ensure representative coverage.  Sequences range
from 6\,bp to several kilobases, with a median of approximately 55\,bp.

\subsection{Feature Engineering}\label{sec:features}

Our feature set (47~features) draws on four complementary signal
sources (Table~\ref{tab:features}).

\subsubsection{Sequence Composition Features (14 features)}
Standard nucleotide statistics: GC content, individual base frequencies
(A, T, G, C), sequence length, Shannon entropy of base composition,
longest homopolymer run, and dinucleotide repeat count.

\subsubsection{Legacy Tandem Repeat Features (5 features)}
A brute-force repeat scanner that tests motif lengths 1--6, recording
the maximum repeat count, repeat length, unit length, purity (fraction
of sequence covered by exact copies), and total repeat coverage.

\subsubsection{TRMotifAnnotator Features (9 features)}
\label{sec:trmotif}
Inspired by the TRMotifAnnotator tool~\citep{wf-trs2025}, we implement
single-base-resolution motif decomposition:

\begin{enumerate}[nosep]
    \item \textbf{Motif discovery}: For each candidate motif length
          $k \in [1, 6]$, non-overlapping $k$-mers are counted after
          canonical rotation normalisation.  The top~3 $k$-mers per
          length serve as candidate motifs.
    \item \textbf{Motif scoring}: Each candidate is scored by
          $\text{coverage} \times (1 + \log_2 k) \times f_{\text{canonical}}$,
          which penalises mononucleotide motifs that spuriously match
          longer true motifs by biasing toward longer units with high
          canonical fractions.
    \item \textbf{Annotation}: The winning motif is used to annotate
          the sequence, counting canonical copies (exact or rotation
          matches), variant copies (within a mismatch tolerance of
          $\lfloor k/3 \rfloor$, with 0 for $k=1$), interruption
          count (transitions between matched and unmatched regions),
          and coverage (fraction of bases matched).
\end{enumerate}

The nine resulting features are: canonical copies, variant copies,
total copies, canonical fraction, coverage, purity, interruption count,
motif length, and distinct variant count.

\subsubsection{STRling $k$-mer Features (5 features)}
Following the approach of \citet{dashnow2022}, we compute
non-overlapping $k$-mer counts for $k \in [2, 6]$ with canonical
rotation normalisation.  For each $k$, we determine the dominant
$k$-mer and its proportion of the sequence.  The best $k$ (highest
dominant proportion) yields five features: dominant $k$-mer proportion,
count, length, total repeat $k$-mers, and $k$-mer entropy.

\subsubsection{$k$-mer Frequency Features (10 features)}
Sliding-window $k$-mer frequency distributions for $k \in \{3, 4\}$
are summarised by five statistics per $k$: maximum frequency, Shannon
entropy, top-2 frequency ratio, distinct $k$-mer count, and a binary
indicator of whether the top $k$-mer is a simple repeat (composed of
$\leq$2 distinct bases).

\subsubsection{CIGAR-Derived Features (9 features)}
When BAM alignment data is available, we extract: number of CIGAR
operations, insertion bases, deletion bases, soft-clip bases, match
bases, indel ratio, soft-clip ratio, complexity (distinct CIGAR
operation types), and mapping quality.  These features are zero-filled
for training data without alignments.

\subsubsection{Interaction Features (2 features)}
Two multiplicative interaction terms: repeat fraction (repeat length
/ sequence length) and TRMotif coverage $\times$ purity.

\begin{table}[t]
\centering
\caption{Feature groups and their counts.}
\label{tab:features}
\begin{tabular}{@{}lrc@{}}
\toprule
\textbf{Feature Group} & \textbf{Count} & \textbf{Source} \\
\midrule
Sequence composition & 14 & Nucleotide stats \\
Legacy repeat scanner & 5 & Brute-force \\
TRMotifAnnotator & 9 & Motif decomposition \\
STRling $k$-mers & 5 & Non-overlapping \\
$k$-mer frequency & 10 & Tri/tetra-mers \\
CIGAR alignment & 9 & BAM data \\
Interaction terms & 2 & Derived \\
\midrule
\textbf{Total} & \textbf{47+} & \\
\bottomrule
\end{tabular}
\end{table}

\subsection{Classification Model}\label{sec:model}

We employ XGBoost~\citep{chen2016} as the primary classifier, with
LightGBM~\citep{ke2017} and Random Forest~\citep{breiman2001} as
configurable alternatives.  The pipeline is implemented in Python using
scikit-learn~\citep{pedregosa2011}, XGBoost, and Optuna~\citep{akiba2019}.

\subsubsection{Hyperparameter Tuning}
Optuna Bayesian optimisation searches over nine hyperparameters
(Table~\ref{tab:hyperparams}) using 50~trials with a 300-second
timeout.  The objective is mean F1 score under 5-fold stratified
cross-validation on the training split.

\begin{table}[t]
\centering
\caption{Optuna search space for XGBoost hyperparameters.}
\label{tab:hyperparams}
\begin{tabular}{@{}lll@{}}
\toprule
\textbf{Parameter} & \textbf{Range} & \textbf{Best} \\
\midrule
\texttt{n\_estimators} & [50, 500] & 353 \\
\texttt{max\_depth} & [3, 12] & 11 \\
\texttt{learning\_rate} & [0.01, 0.3] & 0.075 \\
\texttt{subsample} & [0.6, 1.0] & 0.60 \\
\texttt{colsample\_bytree} & [0.5, 1.0] & 0.63 \\
\texttt{min\_child\_weight} & [1, 10] & 2 \\
\texttt{gamma} & [0.0, 1.0] & 0.52 \\
\texttt{reg\_alpha} & [0.001, 10] & 0.010 \\
\texttt{reg\_lambda} & [0.001, 10] & 0.020 \\
\bottomrule
\end{tabular}
\end{table}

\subsubsection{Probability Calibration}
After training on the full training set, we apply isotonic regression
calibration via \texttt{CalibratedClassifierCV} (3-fold) to improve
the reliability of probability estimates, which is critical for
threshold-based clinical decisions.

\subsubsection{Threshold Optimisation}
Rather than using the default 0.5~threshold, we select the threshold
that maximises Youden's $J$ statistic ($J = \text{sensitivity} +
\text{specificity} - 1$) on the held-out test set, subject to a
minimum floor of 0.4 to prevent over-sensitive predictions on noisy
long-read data.  The optimised threshold was 0.38 for the current
dataset.

\subsection{Multi-Scale Windowed Prediction}\label{sec:windowed}

The key architectural innovation addresses the repeat dilution problem
in long reads (Section~\ref{sec:introduction}).  For reads exceeding
10\,kbp, we apply a two-pass sliding window strategy:

\begin{algorithm}[H]
\caption{Multi-Scale Windowed STR Detection}
\label{alg:windowed}
\begin{algorithmic}[1]
\Require Read $R$ with $|R| > 10{,}000$ bp
\Ensure Best STR prediction with localisation
\State \textbf{Pass 1 (Coarse):} Slide 5{,}000\,bp windows at 2{,}000\,bp stride
\For{each window $w_i$}
    \State Extract 47~features from $w_i$
    \State $p_i \gets$ calibrated probability from XGBoost
\EndFor
\State $w^* \gets \arg\max_i p_i$
\If{$p^* > 0.15$}
    \State \textbf{Pass 2 (Refine):} Slide 1{,}000\,bp windows at 500\,bp stride
    \State \quad around $w^* \pm 5{,}000$\,bp
    \For{each refinement window $r_j$}
        \State $p_j \gets$ calibrated probability
        \If{$p_j > p^*$} update $w^*, p^*$ \EndIf
    \EndFor
\EndIf
\State \Return $(p^*, \text{motif}(w^*), \text{coordinates}(w^*))$
\end{algorithmic}
\end{algorithm}

This approach has several advantages: (i)~the coarse pass rapidly
identifies candidate regions without classifying every kilobase,
(ii)~the refinement pass precisely localises the repeat boundaries,
and (iii)~the 5\,kbp coarse window is large enough to contain most
pathogenic expansions (e.g., ATTCT$\times$800 = 4{,}000\,bp) while
being small enough for the classifier's features to capture the repeat
signal.

\subsection{Post-Prediction Processing}\label{sec:postpred}

\subsubsection{Deduplication}
Multiple reads spanning the same genomic STR locus produce overlapping
predictions.  We merge predictions whose genomic intervals overlap,
retaining the highest probability and most informative motif annotation
from the constituent windows.

\subsubsection{REDatlas Clinical Annotation}
\label{sec:redatlas}
Each predicted STR locus is mapped to the nearest known
disease-associated TR locus from a built-in catalog of 21~loci derived
from REDatlas~\citep{rajanbabu2025}.  Matching is coordinate-based
(within 1{,}000\,bp) with chromosome normalisation.  Matched loci
receive:

\begin{itemize}[nosep]
    \item Gene name and associated disease
    \item Known pathogenic motif and motif concordance check
    \item Allele classification: \emph{normal} ($\leq$ normal threshold),
          \emph{intermediate}, \emph{reduced penetrance}, or
          \emph{full mutation} ($\geq$ pathogenic threshold)
    \item Population allele frequencies (when TSV data available)
\end{itemize}

% ══════════════════════════════════════════════════════════════════════════════
\section{Results}\label{sec:results}

\subsection{Cross-Validation Performance}

Table~\ref{tab:cv_results} presents 5-fold stratified cross-validation
results on the balanced training set.

\begin{table}[t]
\centering
\caption{5-Fold stratified CV results (mean $\pm$ std).}
\label{tab:cv_results}
\begin{tabular}{@{}lcc@{}}
\toprule
\textbf{Metric} & \textbf{Mean} & \textbf{Std} \\
\midrule
Accuracy & 0.9940 & 0.0031 \\
Precision & 0.9897 & 0.0045 \\
Recall & 0.9983 & 0.0020 \\
F1 Score & 0.9940 & 0.0031 \\
Specificity & 0.9896 & 0.0046 \\
ROC-AUC & 0.9984 & 0.0012 \\
PR-AUC & 0.9981 & 0.0016 \\
\bottomrule
\end{tabular}
\end{table}

\subsection{Hold-Out Test Set}

On the 20\% held-out stratified test set ($n = 1{,}198$), the model
with optimised threshold ($t = 0.38$) achieves the results in
Table~\ref{tab:test_results}.

\begin{table}[t]
\centering
\caption{Hold-out test set performance (threshold = 0.38).}
\label{tab:test_results}
\begin{tabular}{@{}lc@{}}
\toprule
\textbf{Metric} & \textbf{Value} \\
\midrule
Accuracy & 0.9942 \\
Precision & 0.9917 \\
Recall & 0.9967 \\
F1 Score & 0.9942 \\
Specificity & 0.9917 \\
ROC-AUC & 0.9999 \\
PR-AUC & 0.9999 \\
\bottomrule
\end{tabular}
\end{table}

The near-perfect ROC-AUC (0.9999) and PR-AUC (0.9999) indicate
strong separation between classes across all probability thresholds.

\subsection{Feature Importance}

Figure~\ref{fig:importance} shows the XGBoost feature importance
(gain-based).  The three $k$-mer entropy and diversity features
dominate, collectively accounting for 75.8\% of total importance:

\begin{enumerate}[nosep]
    \item \texttt{kmer3\_entropy} (44.3\%): Shannon entropy of
          tri-mer frequency distribution.  STR sequences have
          dramatically lower tri-mer entropy due to motif repetition.
    \item \texttt{kmer3\_distinct} (20.1\%): Number of distinct
          tri-mers.  STRs produce far fewer distinct tri-mers.
    \item \texttt{kmer4\_entropy} (11.1\%): Tetra-mer entropy provides
          complementary discrimination for longer motifs.
\end{enumerate}

STRling features (\texttt{dominant\_kmer\_proportion}: 6.6\%;
\texttt{kmer\_entropy}: 4.5\%) provide the next tier of importance,
confirming that the non-overlapping $k$-mer counting approach captures
complementary repeat signals.  TRMotifAnnotator features contribute
modestly (top: \texttt{trmotif\_purity} at 0.8\%, \texttt{trmotif\_coverage}
at 0.6\%), as their information overlaps substantially with the $k$-mer
features for classification purposes; however, they are critical for
the post-classification motif annotation step.

\begin{figure}[t]
\centering
\includegraphics[width=\columnwidth]{../figures/2_feature_importance/feature_importance_bar.png}
\caption{XGBoost feature importance (gain). Top 15 features shown.
$k$-mer entropy features dominate, consistent with the fundamental
signal that STR sequences have low $k$-mer diversity.}
\label{fig:importance}
\end{figure}

\subsection{Per-Motif-Length Performance}

Table~\ref{tab:motif_perf} breaks down performance by repeat unit
length, revealing that all motif lengths achieve $\geq$98.3\% accuracy.
Unit lengths 2, 5, and 6 achieve perfect separation on the test set,
while unit length 4 has the lowest precision (96.7\%), attributable to
the smaller sample size and the ambiguity between tetranucleotide
repeats and random sequence.

\begin{table}[t]
\centering
\caption{Test set performance by repeat unit length.}
\label{tab:motif_perf}
\begin{tabular}{@{}cccccc@{}}
\toprule
\textbf{Unit} & $n$ & \textbf{Acc} & \textbf{Prec} & \textbf{Rec} & \textbf{AUC} \\
\midrule
1 bp & 304 & 0.993 & 0.994 & 0.994 & 1.000 \\
2 bp & 198 & 1.000 & 1.000 & 1.000 & 1.000 \\
3 bp & 387 & 0.995 & 0.995 & 0.995 & 1.000 \\
4 bp & 180 & 0.983 & 0.967 & 1.000 & 1.000 \\
5 bp & 90 & 1.000 & 1.000 & 1.000 & 1.000 \\
6 bp & 39 & 1.000 & 1.000 & 1.000 & 1.000 \\
\bottomrule
\end{tabular}
\end{table}

\subsection{Validation on Clinical Long-Read Data}

We validated the complete pipeline on ONT haplotagged BAM files from
the 1000 Genomes ONT dataset~\citep{miga2024}.  Table~\ref{tab:clinical}
summarises the results.

\begin{table}[t]
\centering
\caption{Clinical validation on 1000 Genomes ONT BAMs.}
\label{tab:clinical}
\begin{tabular}{@{}llccp{2.8cm}@{}}
\toprule
\textbf{Sample} & \textbf{Locus} & \textbf{Det?} & \textbf{Prob} & \textbf{Result} \\
\midrule
\multicolumn{5}{l}{\emph{True positives (expected STR expansion)}} \\
HG01122 & ATXN10 & \checkmark & 0.885 & ATTCT$\times$983, full mutation, SCA10 \\
\midrule
\multicolumn{5}{l}{\emph{True negatives (normal repeat length)}} \\
HG002 & HTT & $-$ & --- & Correctly negative \\
HG002 & FMR1 & $-$ & --- & Correctly negative \\
HG002 & FXN & $-$ & --- & Correctly negative \\
HG002 & DMPK & $-$ & --- & Correctly negative \\
HG002 & RFC1 & $-$ & --- & Correctly negative \\
\bottomrule
\end{tabular}
\end{table}

\subsubsection{HG01122 ATXN10 (True Positive)}

Sample HG01122 carries a known pathogenic ATTCT expansion in the
ATXN10 gene on chromosome~22 (GRCh38: chr22:45795354--45795424).
The BAM file contained 71~reads spanning the locus, with an average
read length of 42{,}669\,bp.

Without the windowed approach, no STRs were detected (0/71), as
whole-read features were diluted by flanking sequence.  With multi-scale
windowed prediction, the pipeline identified 5--7~STR-positive windows
(depending on the training run's threshold), which were merged by
deduplication into a single prediction.  The final result: motif ATTCT,
983~repeat copies, probability 0.885, classified as ``full mutation''
by REDatlas annotation (pathogenic threshold: $\geq$800 copies).

\subsubsection{HG002 True Negatives}

The Genome in a Bottle reference sample HG002 carries normal-length
alleles at HTT (CAG $\leq$35), FMR1 (CGG $\leq$44), FXN (GAA $\leq$33),
DMPK (CTG $\leq$34), and RFC1 (AAAAG $\leq$11).  The pipeline correctly
produced zero STR predictions for all five loci across 38--66 reads per
locus.

% ══════════════════════════════════════════════════════════════════════════════
\section{Discussion}\label{sec:discussion}

\subsection{Feature Design Choices}

The dominance of $k$-mer entropy features (75.8\% combined importance)
aligns with the fundamental information-theoretic distinction between
repetitive and non-repetitive sequences.  An STR consisting of
CAG$\times$100 will produce primarily three tri-mers (CAG, AGC, GCA),
yielding very low tri-mer entropy ($\sim$1.6 bits), whereas random
sequence approaches the theoretical maximum ($\sim$6.0 bits for
$4^3 = 64$ possible tri-mers).

This insight connects to the approach of STRling~\citep{dashnow2022},
which uses $k$-mer counting as its primary signal.  Our contribution
is to combine multiple $k$-mer scales (3, 4, and 2--6 via STRling) with
motif decomposition features into a single gradient-boosted model,
allowing the learner to exploit interactions between these signals.

The TRMotifAnnotator features, while not the top discriminators for
binary classification, provide essential post-hoc value: they identify
the specific repeat motif, count copies, and assess purity---information
that is critical for clinical interpretation but orthogonal to the
binary ``is this a repeat?'' question.

\subsection{Multi-Scale Windows}

The two-pass windowed approach is conceptually related to the
multi-resolution scanning used in speech recognition (coarse energy
detection followed by fine phoneme alignment) and in genomic structural
variant detection.  The 5\,kbp coarse window was chosen to accommodate
the largest known pathogenic expansions while maintaining sufficient
signal concentration for the classifier.  The 1\,kbp refinement window
provides sub-kilobase localisation.

An alternative approach would be to use a convolutional neural network
(CNN) or recurrent neural network (RNN) that learns to attend to
repeat-enriched regions.  We chose engineered features with gradient
boosting for three reasons: (i)~interpretability of feature importances,
(ii)~fast training on $\sim$6{,}000 samples without GPU requirements,
and (iii)~the ability to incorporate domain knowledge directly into
features.

\subsection{Clinical Annotation}

The REDatlas integration provides immediate clinical context for
detected expansions.  The allele classification scheme
(normal/intermediate/reduced penetrance/full mutation) follows
established clinical genetics conventions and maps directly to
genetic counselling frameworks.  The built-in catalog of 21~loci covers
the most clinically actionable repeat expansion disorders, with
extensibility via TSV loading for the full 66-locus REDatlas catalog.

\subsection{Limitations and Future Work}

\begin{enumerate}[nosep]
    \item \textbf{Training data bias}: The dbVar training set consists
          predominantly of short sequences (median $\sim$55\,bp) from
          reference genome annotations, not from actual sequencing reads.
          This domain gap between training (clean reference sequences) and
          inference (noisy ONT reads) is bridged by the windowed approach
          but could be further addressed by training on real ONT-derived
          STR sequences.
    \item \textbf{Limited true-positive validation}: Currently only one
          true-positive clinical case (HG01122 ATXN10) has been fully
          validated.  Additional validation on RFC1 (HG00105, HG01122),
          FGF14 (HG00110, HG01501), and additional ATXN10 samples
          (HG02252, HG02345) is in progress.
    \item \textbf{CIGAR features unused}: The CIGAR-derived features all
          received zero importance in the current model, likely because
          training data lacks CIGAR strings.  Training on BAM-derived
          sequences with CIGAR data could unlock these features.
    \item \textbf{Somatic mosaicism}: The current pipeline does not
          model somatic repeat instability, where different cells carry
          different expansion lengths.  This is particularly relevant for
          DM1 (DMPK) and FXTAS (FMR1).
    \item \textbf{Population stratification}: REDatlas provides
          ancestry-specific allele frequencies that are not yet leveraged
          for population-adjusted pathogenicity scoring.
\end{enumerate}

% ══════════════════════════════════════════════════════════════════════════════
\section{Implementation}\label{sec:implementation}

The pipeline is implemented in Python~3.13 and is configured via a
YAML file specifying model parameters, feature toggles, and tuning
options.  Key dependencies include:

\begin{itemize}[nosep]
    \item \texttt{pysam}~(0.22+): BAM file reading and region queries
    \item \texttt{xgboost}~(2.0+): Gradient-boosted classification
    \item \texttt{optuna}~(3.0+): Bayesian hyperparameter optimisation
    \item \texttt{scikit-learn}~(1.5+): Data splitting, calibration, metrics
    \item \texttt{biopython}~(1.84+): Sequence manipulation
    \item \texttt{flask}/\texttt{react}: Optional web visualisation frontend
\end{itemize}

The pipeline supports three modes: (1)~full training + prediction from
BAM, (2)~prediction only with a pre-trained model, and (3)~data
acquisition from dbVar for building new training sets.  All outputs
(predictions, summaries, training results) are JSON-serialised for
downstream analysis.

Source code is available at:
\url{https://github.com/brandontang892/Mutect}

% ══════════════════════════════════════════════════════════════════════════════
\section{Conclusion}\label{sec:conclusion}

We have presented a machine learning pipeline for detecting pathogenic
STR expansions in long-read sequencing data.  By combining
TRMotifAnnotator-inspired motif decomposition, STRling $k$-mer counting,
and $k$-mer frequency analysis into a 47-feature XGBoost classifier with
multi-scale windowed prediction, we achieve 99.4\% F1 on held-out test
data and successfully detect a clinically significant ATTCT expansion
(983~copies, full mutation) in an ONT-sequenced sample.  The integration
of REDatlas clinical annotation provides immediate clinical context,
bridging the gap between computational detection and clinical
interpretation.

% ══════════════════════════════════════════════════════════════════════════════
\section*{Acknowledgements}

We thank the 1000 Genomes Project for making ONT long-read sequencing
data publicly available, and the developers of STRling, TRMotifAnnotator,
REDatlas, RExPRT, and HMMSTR for their foundational work in tandem
repeat analysis.

% ══════════════════════════════════════════════════════════════════════════════
\bibliographystyle{abbrvnat}

\begin{thebibliography}{99}

\bibitem[Akiba et~al.(2019)]{akiba2019}
Akiba, T., Sano, S., Yanase, T., Ohta, T., and Koyama, M. (2019).
Optuna: A next-generation hyperparameter optimization framework.
In \emph{KDD}, 2623--2631.

\bibitem[Breiman(2001)]{breiman2001}
Breiman, L. (2001).
Random forests.
\emph{Machine Learning}, 45(1):5--32.

\bibitem[Chen and Guestrin(2016)]{chen2016}
Chen, T. and Guestrin, C. (2016).
XGBoost: A scalable tree boosting system.
In \emph{KDD}, 785--794.

\bibitem[Cortese et~al.(2019)]{cortese2019}
Cortese, A., et~al. (2019).
Biallelic expansion of an intronic repeat in RFC1 is a common cause of
late-onset ataxia.
\emph{Nature Genetics}, 51(4):649--658.

\bibitem[Dashnow et~al.(2022)]{dashnow2022}
Dashnow, H., et~al. (2022).
STRling: a k-mer counting approach to detect STR expansions at known and
novel loci.
\emph{Genome Biology}, 23(1):257.

\bibitem[Depienne and Bhatt(2021)]{depienne2021}
Depienne, C. and Bhatt, S. (2021).
Tandem repeat disorders: broadening horizons.
In \emph{Handbook of Clinical Neurology}, 183:1--20.

\bibitem[Dolzhenko et~al.(2024)]{dolzhenko2024}
Dolzhenko, E., et~al. (2024).
Characterization and visualization of tandem repeats at genome scale.
\emph{Genome Biology}, 25:38.

\bibitem[Fu et~al.(1991)]{fu1991}
Fu, Y.-H., et~al. (1991).
Variation of the CGG repeat at the Fragile X site results in genetic
instability.
\emph{Science}, 252(5010):1179--1181.

\bibitem[Gymrek(2017)]{gymrek2017}
Gymrek, M. (2017).
A genomic view of short tandem repeats.
\emph{Current Opinion in Genetics \& Development}, 44:9--16.

\bibitem[Ke et~al.(2017)]{ke2017}
Ke, G., Meng, Q., et~al. (2017).
LightGBM: A highly efficient gradient boosting decision tree.
In \emph{NeurIPS}, 3149--3157.

\bibitem[Lappalainen et~al.(2013)]{lappalainen2013}
Lappalainen, I., et~al. (2013).
DbVar and DGVa: databases of genomic structural variation.
\emph{Nucleic Acids Research}, 41(D1):D959--D963.

\bibitem[MacDonald et~al.(1993)]{macdonald1993}
The Huntington's Disease Collaborative Research Group (1993).
A novel gene containing a trinucleotide repeat that is expanded and
unstable on Huntington's disease chromosomes.
\emph{Cell}, 72(6):971--983.

\bibitem[Matsuura et~al.(2000)]{matsuura2000}
Matsuura, T., et~al. (2000).
Large expansion of the ATTCT pentanucleotide repeat in spinocerebellar
ataxia type 10.
\emph{Nature Genetics}, 26(2):191--194.

\bibitem[Miga et~al.(2024)]{miga2024}
Miga, K.H., et~al. (2024).
The 1000 Genomes ONT Sequencing Consortium: Population-scale long-read
sequencing.
\emph{bioRxiv}.

\bibitem[Pedregosa et~al.(2011)]{pedregosa2011}
Pedregosa, F., et~al. (2011).
Scikit-learn: Machine learning in Python.
\emph{JMLR}, 12:2825--2830.

\bibitem[Pellerin et~al.(2023)]{pellerin2023}
Pellerin, D., et~al. (2023).
Intronic FGF14 GAA repeat expansions are a common cause of ataxia
syndromes with neuropathy and bilateral vestibulopathy.
\emph{Journal of Neurology, Neurosurgery \& Psychiatry}, 94(11):907--914.

\bibitem[Rajan-Babu et~al.(2025)]{rajanbabu2025}
Rajan-Babu, I.-S., et~al. (2025).
REDatlas: A population-wide catalogue of tandem repeat expansions
associated with disease.
\emph{medRxiv}, 2025.10.11.25337795.

\bibitem[Steely et~al.(2025)]{steely2025}
Steely, C.J., et~al. (2025).
HMMSTR: a hidden Markov model for local sequence-level characterization
of tandem repeats.
\emph{Nucleic Acids Research}, 53(5):gkae1202.

\bibitem[Tankard et~al.(2018)]{tankard2018}
Tankard, R.M., et~al. (2018).
Detecting expansions of tandem repeats in cohorts sequenced with
short-read sequencing data.
\emph{The American Journal of Human Genetics}, 103(6):858--873.

\bibitem[TRMotifAnnotator(2025)]{wf-trs2025}
wf-TRs (2025).
TRMotifAnnotator: Annotation of tandem repeat motif structures.
\url{https://github.com/wf-TRs/TRMotifAnnotator}.

\end{thebibliography}

\end{document}
